%----------------------------------------------------------------------------------------
%	PACKAGES AND OTHER DOCUMENT CONFIGURATIONS
%----------------------------------------------------------------------------------------

% Compile this document with XeLaTeX or LuaLaTeX
\documentclass[11pt,a4paper,sans]{moderncv} 

% Modern CV themes
\moderncvstyle{classic} % Style options: 'casual' (default), 'classic', 'banking', 'oldstyle', 'fancy'
\moderncvcolor{blue} % Color options: 'blue' (default), 'orange', 'green', 'red', 'purple', 'grey', 'black'

% Character encoding
\usepackage[utf8]{inputenc}
\usepackage[scale=0.85]{geometry} % Adjust page margins
\usepackage{makecell}
\usepackage[unicode]{hyperref}
\newcommand{\specialcell}[2][c]{%
  \begin{tabular}[#1]{@{}l@{}}#2\end{tabular}}


% Personal data
\name{Saurav}{Bhattarai}
\title{Water Resource Engineer \& Hydro-climatology Researcher}

\phone[mobile]{+1-601-918-0712} % Replace with your actual phone number
\email{Saurav.bhattarai@students.jsums.edu}
\email{saurav.bhattarai.1999@gmail.com}
\social[linkedin]{saurav-bhattarai-7133a3176}
\social[github]{Sauravbhattarai19}
\homepage{sauravbhattarai19.github.io/}
\photo[80pt][0pt]{Headshot.png} % Replace with your photo if desired

%----------------------------------------------------------------------------------------

\begin{document}

\makecvtitle

%----------------------------------------------------------------------------------------
%	OBJECTIVE
%----------------------------------------------------------------------------------------

\section{Objective}
To become a skilled water resource engineer with expertise in hydro-climatology, leveraging advanced programming and modeling tools to contribute to the global challenge of water-resource management, focusing on the impacts of climate change and water-related hazards on society.

%----------------------------------------------------------------------------------------
%	CURRENT POSITIONS
%----------------------------------------------------------------------------------------

\section{Current Positions}

\cventry{2023--Present}{Graduate Research Assistant}{Jackson State University}{Department of Civil and Environmental Engineering}{1400 John R. Lynch St, Jackson, MS 39217}{}

\cventry{2024--Present}{ORISE Graduate Fellow}{U.S. Army Corps of Engineers - Engineer Research and Development Center}{Coastal and Hydraulics Laboratory (ERDC-CHL)}{3909 Halls Ferry Road, Vicksburg, MS 39180}{}

%----------------------------------------------------------------------------------------
%	EDUCATION
%----------------------------------------------------------------------------------------

\section{Education}

\cventry{2023--Present}{Graduate Studies}{Jackson State University}{USA}{GPA: 4.0/4.0}{}

\cventry{2017--2022}{Bachelor of Civil Engineering}{Tribhuvan University, Institute of Engineering (IOE)}{Pulchowk Campus, Nepal}{Score: 80.14\%}{}

\cventry{2015--2017}{Higher Secondary Education}{Capital Secondary School (CCRC)}{Koteshwor, Kathmandu, Nepal}{Score: 85.90\%}{}

\cventry{2015}{Secondary Education}{Karunanidhi Education Foundation School}{Simpani, Pokhara, Nepal}{Score: 91.00\%}{}

%----------------------------------------------------------------------------------------
%	ACHIEVEMENTS
%----------------------------------------------------------------------------------------

\section{Achievements}

\cvitem{2025}{Lead Developer, Microsoft × AUC Data Science Initiative Data Challenge — Led the development of the GeoClimate Platform, recognized as “Best in Show” and awarded First Position among finalists in Atlanta, GA, presented by Dr. Rocky Talchabhadel (December 11, 2025).}

\cvitem{2025}{Team Leader, Mastercard × AUC Data Science Initiative Data Challenge — Led Team "GeoData Catalyst" from Jackson State University to the final showcase in Atlanta, GA; received an Honorable Mention among top 18 teams selected from 70+ participating HBCUs and was awarded \$1,000 (December 4–6, 2025).}

\cvitem{2025}{Selected Participant, CUAHSI-SINTER Snow Measurement Field School, Granby, CO (January 5-8, 2025) -- Competitively selected for intensive training in snow measurement techniques including depth, density, water equivalence, grain analysis, and stratigraphy. Awarded \$600 travel reimbursement.}

\cvitem{2025}{Travel Grant Recipient, Tethys Summit 2025, Tampa, FL -- Selected to attend and present at the annual summit of the Tethys Platform and GEOGLOWS open-source hydrologic ecosystems.}

\cvitem{2025}{Selected Participant, NCAR Graduate Visitor Program (GVP), Boulder, CO -- Chosen to collaborate with scientists at the National Center for Atmospheric Research on climate-hydrology research as part of a competitive summer research residency.}

\cvitem{2024--2025}{ORISE Fellow, U.S. Department of Defense Research Participation Program (ERDC CHL): Selected for a prestigious fellowship focusing on hydrological modeling and analysis.}

\cvitem{2024}{Winner, AGU Catchment Hydrology Subcommittee Early Career Competition, American Geophysical Union -- Awarded \$300 for submission on rainfall-runoff modeling using LSTM sequence-to-sequence learning models.}

\cvitem{2024}{CyberTraining Workshop Travel Support: Awarded a travel grant to attend the CyberTraining workshop organized by the Natural Hazards Center.}

\cvitem{2024}{AGU-Chapman Travel Grant to support attendance at Chapman Conference on Remote Sensing of the Water Cycle, Honolulu, HI, Feb 13-16, 2024.}

\cvitem{2023--Present}{Full Scholarship with a position of Graduate Research Assistant in the Department of Civil and Environmental Engineering under Prof. Dr. Rocky Talchabhadel in Jackson State University from the "Hydrological Impacts Computing, Outreach, and Resiliency Partnership (HICORPS) Project" in collaboration with ERDC and WOOLPERT.}

\cvitem{2017--2021}{Golden Jubilee Scholarship Scheme (GJSS) by H.E. Ambassador of India, Nepal.}

\cvitem{2017--2021}{Merit-Based Full Four-Year Scholarship awarded by Institute of Engineering (IOE), Pulchowk Campus for undergraduate study in Civil Engineering for securing 19\textsuperscript{th} rank out of approximately 15,000 engineering aspirants.}

\cvitem{2015--2017}{Merit-Based Full Two-Year Scholarship awarded by Capital Secondary School (CCRC) for the higher secondary course.}

\cvitem{2014}{Second Position in Western Regional Inter-Secondary School Mathematics Quiz Contest organized by Council for Mathematics Education, Kaski.}

\cvitem{2014}{Served as President of Junior Red Cross Society at KEF High School, providing leadership and oversight of the chapter's activities and community service initiatives.}

\cvitem{2013}{Honored with Certificate of Participation in i-Plea; 1\textsuperscript{st} International Environment Olympiad organized by City Montessori School (CMS), Lucknow, India.}

\cvitem{2012}{First Position in International Mathematics Competition Nepal organized by Asia Association of Education and Exchange (AAEE), Tokyo, Hachioji-shi 40.}

%----------------------------------------------------------------------------------------
%	PROFESSIONAL ORGANIZATION MEMBERSHIP
%----------------------------------------------------------------------------------------

\section{Professional Organization Membership}

\cvitemwithcomment{Member}{American Geophysical Union (AGU)}{}
\cvitemwithcomment{Member}{American Society of Civil Engineers (ASCE)}{}
\cvitemwithcomment{Member}{European Geosciences Union (EGU)}{}
\cvitemwithcomment{Member}{Nepal Red Cross Society (NCRS)}{}
\cvitemwithcomment{Member}{Nepal Engineering Council}{}

%----------------------------------------------------------------------------------------
%	PROJECT/FELLOWSHIP PROPOSALS
%----------------------------------------------------------------------------------------

\section{Project/Fellowship Proposals}

\subsection{Awarded}

\cvitem{2024--2025}{Co-I in Microsoft-AUC Mini-Grants program: "Climate Change Impact Assessment Training for Undergraduate Students: Leveraging Planetary-Scale Data and Cloud"}

\cvitem{2024--2025}{ORISE Fellow, U.S. Department of Defense Research Participation Program (ERDC CHL): Selected for a hydrological modeling and analysis fellowship.}

%----------------------------------------------------------------------------------------
%	PUBLICATIONS
%----------------------------------------------------------------------------------------

%----------------------------------------------------------------------------------------
%	PUBLICATIONS
%----------------------------------------------------------------------------------------

\section{Publications}

\subsection{Published/Accepted}

\cvitem{2025}{Bhattarai, S., Pradhan, N. R., Talchabhadel, R. GeoClimate intelligence platform: A web-based framework for environmental data analysis. Environmental Modelling \& Software, 106826 (2025). https://doi.org/10.1016/j.envsoft.2025.106826}

\cvitem{2025}{Bhattarai, S., He, C., Sharma, S., Kumar, S., Talchabhadel, R. Rainfall Relief or Humidity Havoc? The Boon and Curse of Precipitation During Heatwaves.* Accepted for publication in *Natural Hazards* (2025).}

\cvitem{2025}{Bhattarai, S., Banjara, P., Pandey, V. P., Aryal, A., Pradhan, P., Al-Douri, F., Pradhan, N. R., Talchabhadel, R. Quantifying the cooling effects of blue-green spaces across urban landscapes: A case study of Kathmandu Valley, Nepal. Urban Climate 61, 102493 (2025). https://doi.org/10.1016/j.uclim.2025.102493}

\cvitem{2025}{Talchabhadel, R., Panthi, J., Pandey, V.P., Rakhal, B., Ghimire, G.R., Bista, S., Bhattarai, S., Poudel, S., Bhattarai, Y., Prajapati, R., Thapa, B.R., Nepal, B., Sharma, S. Yesterday's extremes, today's new normal: flood risk in the Kathmandu Valley, Nepal. Natural Hazards 121, 19409–19423 (2025). https://doi.org/10.1007/s11069-025-07524-5}

\cvitem{2025}{Bhattarai, S., Bokati, L., Sharma, S., Talchabhadel, R. Understanding spatiotemporal variation of heatwave projections across US cities. Scientific Reports 15, 10643 (2025). https://doi.org/10.1038/s41598-025-95097-5}

\cvitem{2025}{Bista, S., Bhattarai, S., Sharma, S., Talchabhadel, R. Local-scale flood hazard projections in historically vulnerable communities. Modeling Earth Systems and Environment (Accepted, 2025). https://doi.org/10.1007/s40808-025-02599-2}

\cvitem{2024}{Talchabhadel, R., Bhattarai, S., Bista, S. Projected Changes in Precipitation Extremes Across the Mississippi River Basin Using the NASA Global Daily Downscaled Datasets NEX‐GDDP‐CMIP6. International Journal of Climatology (2024). https://doi.org/10.1002/joc.8748}

\cvitem{2024}{Bhattarai, S., Bista, S., Sharma, S., White, L.D., Amini, F., Talchabhadel, R. Spatiotemporal Characterization of Heatwave Exposure across Historically Vulnerable Communities. Scientific Reports (2024). https://doi.org/10.1038/s41598-024-71704-9}

\cvitem{2024}{Bhattarai, S., Talchabhadel, R. Comparative Analysis of Satellite-Based Precipitation Data across the CONUS and Hawaii: Identifying Optimal Satellite Performance. Remote Sensing 16, 3058 (2024). https://doi.org/10.3390/rs16163058}

\cvitem{2024}{Banjara, P., Bhattarai, S., Pandey, V.P., Talchabhadel, R. Spatiotemporal characterization of heatwaves on an urban center using satellite-based estimates. Theoretical and Applied Climatology (2024). https://doi.org/10.1007/s00704-024-05026-1}

\cvitem{2022}{Bhattarai, T.N., Ghimire, S., Aryal, S., Baaniya, Y., Bhattarai, S., Sharma, S., Bhattarai, P.K., Pandey, V.P. Projected changes in hydro-climatic extremes with CMIP6 climate model outputs: a case of rain-fed river systems in Western Nepal. Stochastic Environmental Research and Risk Assessment (2022). https://doi.org/10.1007/s00477-022-02312-0}

\subsection{Under Review}

\cvitem{2025}{Bhattarai, S., Pradhan, N.R., Talchabhadel, R. Do the Wettest Days Occur Together? A Global Analysis on Disentangling Precipitation Intensity from Seasonal Timing. [Journal of Geophysical Research: Atmospheres]}

\cvitem{2025}{Poudel, S., Bista, S., Bhattarai, S., Sharma, S., Talchabhadel, R. Understanding multi-hazard risk for U.S. coastal cities. [Geomatics, Natural Hazards and Risk]}

\subsection{Preprints}

\cvitem{2025}{Bhattarai, S., Pradhan, N.R., Talchabhadel, R. GeoClimate Intelligence Platform: A Web-Based Framework for Environmental Data Analysis. SSRN Preprint (2025). https://doi.org/10.2139/ssrn.5582694}

\cvitem{2024}{Bhattarai, S., Bista, S., Sharma, S., White, L.D., Amini, F., Talchabhadel, R. Heatwaves in Mississippi: Exploring Variability and Trends in a Changing Climate. SSRN Preprint (2024). https://doi.org/10.2139/ssrn.4698297}


\section{Invited Talks}

\subsection{2024}
\cvitem{}{Thapa, B. R., Talchabhadel, R., and Bhattarai, S. (2024). \textit{Translating Multi-Hazard Risk Assessment to Compounding and Cascading Disaster Risk Management for Resilient Infrastructure}. Invited presentation at the 2nd GeoMandu International Conference, Kathmandu, Nepal, November 28–29.}
%----------------------------------------------------------------------------------------
%	CONFERENCE PAPERS/ABSTRACTS
%----------------------------------------------------------------------------------------

\section{Conference Papers/Abstracts}

\subsection{2025}

\cvitem{}{Gautam, A., Bhattarai, S., Talchabhadel, R.: Unraveling Urban Heat and Moisture Anomalies Across the Continental U.S. Using Local Climate Zones and Earth Observation Data, AGU Fall Meeting 2025, New Orleans, LA, USA, Dec 18, 2025 [Poster Presentation]}

\cvitem{}{Dahal, B., Bhattarai, S., Bhattarai, P. K., Talchabhadel, R.: Evaluating the Applicability of SWOT and SWORD for Nepal’s Water Monitoring: Building a Scalable National Database, AGU Fall Meeting 2025, New Orleans, LA, USA, Dec 16, 2025 [Poster Presentation]}

\cvitem{}{Bhattarai, S., Talchabhadel, R., Pradhan, N. R.: Precipitation Temporal Disaggregation for Event-Based Hydrological Modeling in the Kathmandu Watershed, Nepal, AGU Fall Meeting 2025, New Orleans, LA, USA, Dec 16, 2025 [Poster Presentation]}

\cvitem{}{Talchabhadel, R., Bhattarai, S., Prajapati, A. K., Karn, B., Bhattarai, B., Yadav, R., Bhattarai, P. K.: Spatiotemporal Assessment of Urban Heat Island and Heatwave Dynamics in Densely Populated Nepalese Cities: Evaluating Nature-Based Cooling Strategies, AGU Fall Meeting 2025, New Orleans, LA, USA, Dec 16, 2025 [Poster Presentation]}

\cvitem{}{Pradhan, N. R., Bhattarai, S., Kloss, G., Wagner, A. M., Talchabhadel, R., Deeb, E. J., Ryder, J. L., Tavakoly, A., Courville, Z., Lakshmi, V., Paquette, M., Khona, D., Lewis, M., Carravone, G.: Global Snow Attribute Mapping Tool (Global-SAM), AGU Fall Meeting 2025, New Orleans, LA, USA, Dec 18, 2025 [Oral Presentation]}

\cvitem{}{Talchabhadel, R., Bhattarai, S., Bista, S., Poudel, S., Jones, D., Terrett, S.: Democratizing Earth Science Data for Community-Based Water Hazard Assessment: A Co-Production Platform for Hydrologic Intelligence, AGU Fall Meeting 2025, New Orleans, LA, USA, Dec 18, 2025 [Oral Presentation]}

\cvitem{}{Bhattarai, S., Pradhan, N.R., Talchabhadel, R.: Advancing Groundwater Monitoring and Downscaling Using Open Datasets and Machine Learning in the Mississippi River Basin, XIX World Water Congress, Marrakech, Morocco, Dec 1-5, 2025 [Oral Presentation]}

\cvitem{}{Bhattarai, S., Talchabhadel, R.: Evaluating the Compound Effects of Precipitation and Storm Surge on Coastal Flooding Risk in the CONUS, ASCE EWRI World Environmental \& Water Resources Congress, Anchorage, Alaska, May 20, 2025 [Oral Presentation]}

\cvitem{}{Talchabhadel, R., Bhattarai, S., Bista, S., Poudel, S., Kafoury, R.: Projecting Extreme Temperatures and Precipitation for the Gulf Coast Region of the United States, ASCE EWRI World Environmental \& Water Resources Congress, Anchorage, Alaska, May 20, 2025 [Oral Presentation]}

\subsection{2024}

\cvitem{}{Bhattarai, S., Byrd, A., Talchabhadel, R.: Evaluating the Compound Effects of Precipitation and Storm Surge on Coastal Flooding Risk in the CONUS: A Detailed Case Study of Galveston, AGU Fall Meeting, Washington DC, Dec 9-13, 2024 [Poster Presentation]}

\cvitem{}{Bhattarai, S., Sharma, S., Talchabhadel, R.: Detecting Multihazards in the United States Using Machine Learning: A Comparative Analysis of Individual and Compound Hazard Events, AGU Fall Meeting, Washington DC, Dec 9-13, 2024 [Poster Presentation]}

\cvitem{}{Talchabhadel, R., Bhattarai, S., Banjara, P., and Pandey, V.P., Urban Cooling Strategies: Evaluating the Effectiveness of Green Spaces and Water Bodies in Mitigating Heat Exposure in Kathmandu Valley, Nepal, AGU Fall Meeting, Washington DC, Dec 9-13, 2024 [Poster Presentation]}

\cvitem{}{Bista, S., Bhattarai, S., Sharma, S., and Talchabhadel, R., Implication of non-stationary climate dynamics for flood risk projections. AGU Fall Meeting, Washington, DC, Dec 3-31, 2024 [Online Poster Presentation]}

\cvitem{}{Talchabhadel, R., Bernard, D., Johnson, K., Bhattarai, S., Bista, S., Poudel, S., Kafoury, R.: Evaluation of Climate Change Projections Using Open-Source Software Infrastructure across the Gulf Coast of the United States [NH13B-2293], AGU Fall Meeting, Washington DC, Dec 9-13, 2024 [Poster Presentation]}

\cvitem{}{Bista, S., Bhattarai, S., Sharma, S., Talchabhadel, R.: Implications of non-stationary climate dynamics for flood risk projections [H05-13], AGU Fall Meeting, Washington DC, Dec 9-13, 2024 [Oral Presentation]}

\cvitem{}{Banjara, P., Bhattarai, S., Pandey, V.P., Talchabhadel, R.: Assessing Heatwave Dynamics and Vulnerability in Nepal: A Geospatial Analysis [H34F-04], AGU Fall Meeting, Washington DC, Dec 9-13, 2024 [Oral Presentation]}

\cvitem{}{Talchabhadel, R., Bhattarai, S., Banjara, P., Pandey, V.P.: Urban Cooling Strategies: Evaluating the Effectiveness of Green Spaces and Water Bodies in Mitigating Heat Exposure in Kathmandu Valley, Nepal [GC51H-0059], AGU Fall Meeting, Washington DC, Dec 9-13, 2024 [Poster Presentation]}

\cvitem{}{Talchabhadel, R., S. Bhattarai (2024). "Interconnected Hazards Data", in Natural Hazards Research Summit 2024: Interconnected Hazards: Retrospective Data Analysis for Climate Resilience in the United States. DesignSafe-CI. https://doi.org/10.17603/ds2-a2vx-sx21}

\cvitem{}{Bhattarai, S., Talchabhadel, R.: Assessing Cascading and Compound Climate Extremes and Associated Vulnerabilities Across the United States; 3rd International Conference on Natural Hazards and Risks in a Changing World: Addressing Compound and Multi-Hazard Risk; 12-14 June, 2024; Amsterdam, Netherlands [Oral Presentation]}

\cvitem{}{Bista, S., Bhattarai, Y., Bhattarai, S., Sharma, S., Talchabhadel, R., Amini, F: Urban Flood Inundation Prediction Under Climate Change, 2024 World Environmental \& Water Resources Congress, Milwaukee, WI, May 19 -- 22, 2024 [Oral Presentation]}

\cvitem{}{Talchabhadel, R., Bhattarai, S., Bista, S., White, L., Amini, F: Providing undergraduate students with training on evaluating climate change projections within the framework of Open-source Software Infrastructure, 2024 World Environmental \& Water Resources Congress, Milwaukee, WI, May 19 -- 22, 2024 [Oral Presentation]}

\cvitem{}{Bhattarai, S., Sharma, S., Talchabhadel, R.: Assessing Multidimensional Climate Extremes and Associated Vulnerabilities Across the United States, EGU General Assembly 2024, Vienna, Austria, 14--19 Apr 2024, EGU24-20600, https://doi.org/10.5194/egusphere-egu24-20600 [Oral Presentation]}

\cvitem{}{Talchabhadel, R., Bhattarai, S., Bista, S.: Empowering Undergraduate Students in Remote Sensing for Climate Change Projection Evaluation Using Open-Source Software Infrastructure, AGU Chapman Remote Sensing of the Water Cycle: Sensors to Science to Society, Honolulu, HI, Feb 13-16, 2024. [Poster Presentation]}

\cvitem{}{Bhattarai, S., Talchabhadel, R.: Heatwaves, Humidity, and the Water Cycle: A Remote Sensing Perspective on Climate Transformation, AGU Chapman Remote Sensing of the Water Cycle: Sensors to Science to Society, Honolulu, HI, Feb 13-16, 2024. [Poster Presentation]}

\cvitem{}{Talchabhadel, R., Bhattarai, S., Banjara, P., Sharma, S.: Projecting heatwave consequences in urban centers of the United States under changing climate, American Meteorological Society's 15th Conference on Environment and Health, 104th AMS Annual Meeting, Baltimore, MD, Jan 28 - Feb 1, 2024, [Oral Presentation]}

\subsection{2023}

\cvitem{}{Bhattarai, S., Bokati, L., Sharma, S., Talchabhadel, R.: Assessing Heatwave Impacts and Enhancing Climate Resilience in Highly Populated US Cities, AGU Fall Meeting, San Francisco, Dec 11-15, 2023 [Oral Presentation]}

\cvitem{}{Bhattarai, S., Bista, S., Talchabhadel, R.: Projecting flood risk under climate change in Pearl River basin, Jackson State University Research Engagement Week, Jackson, MS, Oct 26, 2023 [Poster Presentation]}

%----------------------------------------------------------------------------------------
%	PROFESSIONAL INVOLVEMENT
%----------------------------------------------------------------------------------------

\section{Professional Involvement}

\cvitem{2025}{\textbf{Program Manager \& Facilitator:} Mastercard × AUC Data Science Initiative Regional Pitch Stop, Jackson State University, 2025 — Managed and facilitated the regional pitch competition focused on "Data Roots: Local Stories, Regional Impact," coordinating program logistics, guest speakers (Dr. Ebony Dotson, Mr. Jamal Ware), and judging panel while also participating as a team member delivering data-driven solutions for inclusive community growth.}

\cvitem{2025}{\textbf{Lead Instructor:} Python Training Workshop 2025, Jackson State University, May 28–29, 2025 — Designed and led a hands-on Python programming workshop focused on hydroclimate and geospatial applications for undergraduate and graduate students.}

\cvitem{2024}{\textbf{Student Convener:} NH13B - Coastal Storm Risk Modeling and Analysis via Cloud Computing Poster Session, AGU Fall Meeting, Washington DC, Dec 9–13, 2024.}
\cvitem{}{Primary Convener: Thomas C. Massey (USACE-ERDC)}
\cvitem{}{Convener: Rocky Talchabhadel (Jackson State University)}
\cvitem{}{Chairs: Rocky Talchabhadel (Jackson State University), Chris Bender (Taylor Engineering Inc.)}


%----------------------------------------------------------------------------------------
%	PROFESSIONAL EXPERIENCE
%----------------------------------------------------------------------------------------

\section{Professional Experience}

\cventry{May--Aug 2024}{Summer Intern}{USACE ERDC (United State Army Corps of Engineer: Engineering Research and Development Center)}{Vicksburg, Mississippi, USA}{}{}

\cventry{Fall 2023--Present}{Graduate Research Assistant}{Jackson State University}{Jackson, Mississippi, USA}{Under Prof. Dr. Rocky Talchabhadel}{}

\cventry{Spring 2025}{Teaching Assistant}{Water Resources Engineering (CIV 370)}{Jackson State University}{}{}

\cventry{Spring 2024}{Teaching Assistant}{Fluid Mechanics Lab (CIVL 330)}{Jackson State University}{}{}

\cventry{}{Mathematics and Physics Teacher}{Albright Academy}{Pokhara, Nepal}{}{}

%----------------------------------------------------------------------------------------
%	TEACHING EXPERIENCE
%----------------------------------------------------------------------------------------

\section{Teaching Experience}

\cventry{Fall 2025}{Teaching Assistant}{Advanced Water Resources Engineering (CIV 465)}{Jackson State University}{Under Prof. Dr. Rocky Talchabhadel}
{Responsible for developing and delivering instruction on modern data-driven approaches in water resources engineering, including:
\begin{itemize}
\item \textbf{Curriculum Development:} Designed and implemented modules on Big Data Analytics, Water Data Science, Machine Learning applications, Trend Analysis, Spatio-temporal Mapping, and Climate Projections with Bias Correction
\item \textbf{Interactive Learning Platform:} Created comprehensive web-based teaching materials (\texttt{advancedwater.streamlit.app}) integrating visualizations, programming exercises, and real-world datasets in a unified platform enabling hands-on learning
\item \textbf{Advanced Topics Instruction:} Delivered lectures and hands-on training on GIS Mapping, Intensity-Duration-Frequency (IDF) curves, General Extreme Value (GEV) analysis, hydrological modeling, and data-driven hydrological analysis
\item \textbf{Student Assessment:} Developed assignments, managed semester projects, conducted evaluations, and provided mentorship on data analysis techniques and programming applications in water resources
\item \textbf{Modern Pedagogy:} Integrated contemporary tools (Python, GIS, machine learning) to prepare students for data-intensive careers in water resources engineering
\end{itemize}}

\cventry{Spring 2024, 2025}{Teaching Assistant}{Water Resources Engineering (CIV 370)}{Jackson State University}{Under Prof. Dr. Rocky Talchabhadel}
{Led comprehensive instruction covering fundamental through advanced water resources topics, with full responsibility for teaching, assessment design, and grading:
\begin{itemize}
\item \textbf{Core Topics Instruction:} Delivered lectures on Hydrologic Cycle (precipitation, evaporation, infiltration), Frequency Analysis, Hydrologic Modeling (SCS Method, direct runoff), Rainfall-Runoff Analysis (unit hydrograph, S-hydrograph), Hydrological Routing (streamflow and reservoir routing), and Open Channel Hydraulics
\item \textbf{Advanced Applications:} Taught Flood Frequency Analysis with return periods, Stormwater Hydrology including design and management, and emerging topics such as Hydrologic and Hydraulic Design under changing climate conditions
\item \textbf{Assessment Development:} Created comprehensive assignments, homework problems, and examination materials aligned with learning objectives; evaluated student performance through multiple assessment formats
\item \textbf{Statistical Methods:} Instructed students on probability concepts, common data distributions, and their applications to water resources problems including design storm development
\item \textbf{Practical Applications:} Emphasized real-world engineering applications including Rational Method, sewer design, and stormwater management practices
\end{itemize}}

\cventry{Spring 2024}{Teaching Assistant}{Fluid Mechanics Lab (CIVL 330)}{Jackson State University}{Under Prof. Dr. Rocky Talchabhadel}
{Supervised laboratory sessions providing hands-on experience with fundamental fluid mechanics principles through experimental investigations:
\begin{itemize}
\item \textbf{Laboratory Instruction:} Conducted weekly lab sessions covering Bernoulli's Theorem verification, Jet Impact force measurements, Hydrostatic Force determination on submerged surfaces, Viscosity measurements using different viscometers, Density and Specific Gravity measurements, and Laminar/Turbulent Flow characterization including Reynolds number determination
\item \textbf{Student Guidance:} Assisted students in proper experimental setup, data collection procedures, error analysis, and interpretation of results; ensured laboratory safety protocols were followed
\item \textbf{Assessment:} Evaluated laboratory reports, assessed understanding of theoretical principles through experimental applications, and provided feedback on technical writing and data presentation
\item \textbf{Equipment Management:} Maintained laboratory equipment, troubleshot experimental apparatus issues, and ensured proper calibration of measurement instruments
\end{itemize}}

%----------------------------------------------------------------------------------------
%	SOFTWARE SKILLS
%----------------------------------------------------------------------------------------

\section{Software Skills}

\cvitem{Spatial Tools}{QGIS and Google Earth Engine (GEE)}
\cvitem{Programming Tools}{Python}
\cvitem{AI \& ML}{Machine Learning, Artificial Intelligence, and Deep Learning}
\cvitem{Modeling and Analysis}{HEC-HMS and HEC-RAS (Basics), GSSHA (WMS)}
\cvitem{Design and Drafting}{AutoCAD, REVIT, Civil 3D (Basics)}

%----------------------------------------------------------------------------------------
%	REFERENCES
%----------------------------------------------------------------------------------------

\section{References}

\begin{tabular}{@{}p{0.48\textwidth}p{0.48\textwidth}@{}}
\textbf{Prof. Dr. Rocky Talchabhadel} & \textbf{Dr. Nawa Raj Pradhan} \\
\specialcell{
Assistant Professor \\
Department of Civil \& Environmental Engineering \\
Jackson State University \\
Email: rocky.talchabhadel@jsums.edu \\
Phone: +1-915-710-9264
} 
& 
\specialcell{
Research Hydraulic Engineer \\
Coastal \& Hydraulics Laboratory \\
U.S. Army Engineer Research and Development Center \\
Email: Nawa.Pradhan@erdc.dren.mil \\
Phone: +1-601-529-5502
}
\end{tabular}

\begin{center}
\textit{For more information about our research group and ongoing projects, please visit: \href{https://bit.ly/jsu_water}{bit.ly/jsu\_water}}
\end{center}

\end{document}